% Part: lambda-calculus
% Chapter: lambda-definability
% Section: lambda-definable-recursive

\documentclass[../../../include/open-logic-section]{subfiles}

\begin{document}

\olfileid[pt]{lam}{dfl}{ldr}
\olsection{As funções \usetoken{S}{lambda definable} são recursivas}

Não apenas todas as funções recursivas parciais são !!{lambda definable}; a
recíproca também é verdadeira. Isto é, todas as funções !!{lambda definable} são
recursivas parciais.

\begin{thm}
  \ollabel{thm:lambda-computable} Se uma função parcial $f$ é
  !!{lambda definable}, ela é recursiva parcial.
\end{thm}

\begin{proof}
  Esboçamos apenas a prova. Primeiro, aritmetizamos os termos-$\lambd$,
  O.s.\@, atribuímos sistematicamente números de G\"odel aos termos-$\lambd$, usando
  a codificação usual de sequências por potências de primos. Em seguida, definimos uma
  função recursiva parcial $\fn{normalize}(t)$ que opera sobre o
  número de G\"odel~$t$ de um termo lambda como argumento e devolve
  o número de G\"odel da forma normal, se ele tiver uma, ou permanece indefinida
  caso contrário. Definimos então duas funções recursivas parciais,
  $\fn{toChurch}$ e $\fn{fromChurch}$, que levam números naturais aos
  números de G\"odel do numeral de Church correspondente e vice-versa.

  Usando essas funções recursivas, podemos definir a função~$f$ como uma
  função recursiva parcial. Existe um termo-$\lambd$~$F$ que
  !!{lambda define}s~$f$. Para calcular $f(n_1, \dots, n_k)$, obtemos primeiro
  os números de G\"odel dos numerais de Church correspondentes usando
  $\fn{toChurch}(n_i)$ e os anexamos a $\Gn{F}$ para obter o número de G\"odel
  do termo $F \num{n_1}\dots\num{n_k}$. Aplicamos então
  $\fn{normalize}$ a esse número de G\"odel. Se $f(n_1, \dots, n_k)$ está
  definida, $F \num{n_1}\dots\num{n_k}$ tem uma forma normal (que deve
  ser um numeral de Church); caso contrário, não tem forma normal (e, portanto,
  \[\fn{normalize}(\Gn{F\num{n_1}\dots\num{n_k}})\] fica indefinida).
  Por fim, aplicamos $\fn{fromChurch}$ ao número de G\"odel do
  termo normalizado.
\end{proof}

\end{document}
